\documentclass{article}
\usepackage{amsmath,amssymb}
\begin{document}
\section*{Independent audit of the fourth-order auxiliary-spin bound}
The bound $10k\le N+10$ for an AC subspace of order at least four and
$N\ge9$ is valid.  Here is an independent derivation of the normalized
trace formula, using Cartesian spin contractions instead of summing
total-spin multiplicities.

Write $q=j(j+1)$ and $Q=s(s+1)$, and let $\tau$ be normalized trace on
the physical spin-$j$ space.  Rotational covariance and the commutators
$[J_a,J_b]=i\epsilon_{abc}J_c$ give
\[
 \tau(J_aJ_b)=\frac q3\delta_{ab},\qquad
 \tau(J_aJ_bJ_c)=\frac{iq}{6}\epsilon_{abc}.
\]
For four factors, rotational covariance and cyclicity give
\[
 \tau(J_aJ_bJ_cJ_d)=
 A(\delta_{ab}\delta_{cd}+\delta_{ad}\delta_{bc})+
 B\delta_{ac}\delta_{bd}.
\]
The contractions $\sum_{a,b}\tau(J_a^2J_b^2)=q^2$ and
$\sum_{a,b}\tau(J_aJ_bJ_aJ_b)=q(q-1)$ imply
\[
 12A+3B=q^2,\qquad6A+9B=q(q-1),
 \qquad
 A=\frac{q(2q+1)}{30},\quad B=\frac{q(q-2)}{15}.
\]
The second contraction follows from
$\sum_aJ_aJ_bJ_a=(q-1)J_b$, itself obtained by expanding the double
commutator $\sum_a[J_a,[J_a,J_b]]=2J_b$.

For $X=\sum_a J_a\otimes S_a$, the physical partial traces are scalar:
\begin{align*}
 (\tau\otimes\mathrm{id})(I)&=I,&
 (\tau\otimes\mathrm{id})(X)&=0,\\
 (\tau\otimes\mathrm{id})(X^2)&=\frac{qQ}{3}I,&
 (\tau\otimes\mathrm{id})(X^3)&=-\frac{qQ}{6}I,\\
 (\tau\otimes\mathrm{id})(X^4)&=
 \frac{qQ\bigl((3q-1)Q-q+2\bigr)}{15}I.
\end{align*}
For example, the cubic identity uses
$\sum_{abc}\epsilon_{abc}S_aS_bS_c=iQI$; the quartic one uses
$\sum_{ab}S_aS_bS_aS_b=Q(Q-1)I$.
At $s=9/2$, $Q=99/4$, these reproduce exactly the four moments in
\texttt{spin4-proof.tex}.

Put
\[
 x_r=jr+\frac{r(r+1)-99/4}{2},
 \qquad r=-9/2,-7/2,\ldots,9/2,
\]
and
\[
 F(x)=(x-x_{-9/2})(x-x_{-3/2})(x-x_{-1/2})(x-x_{7/2}).
\]
Applying the Cartesian moment formulas to the coefficients of $F$
gives, independently of any spectral trace sum,
\[
 c=(\tau\otimes\mathrm{id})(F(X))/I
 =\frac{9(j+1)(4j^3+724j^2+2637j+6750)}{20}>0.
\]
The Casimir identity $2X=(J+S)^2-J^2-S^2$ gives the displayed
eigenvalues.  They strictly increase since
$x_{r+1}-x_r=j+r+1>0$ for $j\ge9/2$.
The lower roots include the first node; the two middle roots are
adjacent nodes; the last root is the penultimate node.  Hence
$F(x_r)\le0$ for every $r<9/2$, and so $F(x_r)-c<0$ there.
At the top node,
\[
 F(x_{9/2})-c=
 \frac{9j(596j^3+4972j^2+14489j+11688)}{20}>0.
\]

Fourth-order anticoherence makes the compression of $F(X)-cI$ to
$C\otimes V_{9/2}$ zero.  All the spectral signs above are strict.
Projection from that $10k$-dimensional isotropic subspace onto the
positive spectral space is injective: a nonzero kernel vector would
lie entirely in the negative space and have strictly negative quadratic
form.  The positive spectral space is the top spin-$j+9/2$ sector,
of dimension $N+10$, proving $10k\le N+10$.

The accompanying audit script checks the Cartesian contractions, the
resulting trace polynomial, all spectral signs and the top eigenvalue
symbolically.  No construction or general optimality assertion follows
from this upper bound alone.
\end{document}
